\section{The Leray--Serre spectral sequence and twists}\label{sec:spectral_seq_and_twist}

This section gives a geometric interpretation of the f\/iltration of $H^3_G(X; \Z)$ for the Leray--Serre spectral sequence through types of twists. This is carried out by identifying the Leray--Serre spectral sequence with another natural spectral sequence which computes the Borel equivariant cohomology.

Throughout this section, we assume that $G$ is a f\/inite group acting from the left on a `reasonable' space~$X$, such as a locally contractible, paracompact and regular topological space as in~\cite{FHT}, or a $G$-CW complex~\cite{May}.

\subsection{Spectral sequences}

The Borel equivariant cohomology $H^n_G(X; \Z)$ is def\/ined to be the (singular) cohomology of the quotient space $EG \times_G X$ of $EG \times X$ under the diagonal $G$-action $(\xi, x) \mapsto (\xi g, g^{-1}x)$, where $EG$ is the total space of the universal $G$-bundle $EG \to BG$. Associated to the f\/ibration $X \to EG \times_G X \to BG$ is the Leray--Serre spectral sequence
\begin{gather*}
E_r^{p, q} \Longrightarrow H^{p+q}_G(X; \Z)
\end{gather*}
converging to the graded quotient of a f\/iltration
\begin{gather*}
H^n_G(X; \Z) = F^0H^n_G(X; \Z) \supset F^1H^n_G(X; \Z) \supset \cdots \supset F^{n+1}H^n_G(X; \Z) = 0,
\end{gather*}
that is, $E_\infty^{p, q} = F^pH^{p+q}_G(X; \Z)/F^{p+1}H^{p+q}_G(X; \Z)$. The $E_2$-term is given by the group cohomology of~$G$
\begin{gather*}
E_2^{p, q} = H^p_{\mathrm{group}}(G; H^q(X; \Z)),
\end{gather*}
where the coef\/f\/icient $H^q(X; \Z)$ is regarded as a right $G$-module by the pull-back action. As a~convention of this paper, the group of $p$-cochains with values in a right $G$-module $M$ is denoted by $C^p_{\mathrm{group}}(G; M) = C(G^p, M) = \{ \tau \colon G^p \to M \}$, and the coboundary $\partial \colon C^p_{\mathrm{group}}(G; M) \to C^{p+1}_{\mathrm{group}}(G; M)$ is given by
\begin{gather*}
(\partial \tau)(g_1, \ldots, g_{p+1})= \tau(g_2, \ldots, g_{p+1})+ \sum_{i = 1}^p (-1)^i \tau(g_1, \ldots, g_ig_{i+1}, \ldots, g_{p+1}) \\
\hphantom{(\partial \tau)(g_1, \ldots, g_{p+1})=}{}+ (-1)^{p+1} \tau(g_1, \ldots, g_p)g_{p+1}.
\end{gather*}

As an application of the spectral sequence, we can obtain an identif\/ication $H^n_G(\pt; \Z) \cong H^n_{\mathrm{group}}(G; \Z)$. (We also have $H^n_{\mathrm{group}}(G; \Z) \cong H^{n-1}_{\mathrm{group}}(G; U(1))$ for $n \ge 2$ by the so-called exponential exact sequence.)
